\chapter{مقدمه‌ای بر رگرسیون}
\label{ch:regression}

\section*{اهداف این فصل}
\begin{itemize}
    \item درک تفاوت رگرسیون با طبقه‌بندی
    \item یادگیری رگرسیون خطی ساده و چندگانه
    \item آشنایی با رگرسیون لجستیک به‌عنوان پل میان رگرسیون و طبقه‌بندی
    \item آشنایی با معیارهای ارزیابی مدل‌های رگرسیون
\end{itemize}

\section{رگرسیون در برابر طبقه‌بندی}

در دو فصل قبل، با \textbf{طبقه‌بندی} آشنا شدیم که در آن متغیر هدف
\textbf{کیفی} یا گسسته است (مثلاً «اسپم» یا «غیراسپم»).
\dmterm{\textbf{رگرسیون}}{Regression} نیز مانند طبقه‌بندی یک وظیفهٔ
یادگیری با نظارت است، با این تفاوت اساسی که متغیر هدف در آن
\textbf{عددی و پیوسته} است — مانند پیش‌بینی قیمت یک خانه، دمای فردا، یا
میزان فروش ماه آینده. بسیاری از مفاهیمی که در فصل‌های قبل دیدیم (تفکیک
آموزش/آزمون، اعتبارسنجی متقاطع، خطر بیش‌برازش) بدون تغییر برای رگرسیون
نیز صادق‌اند؛ آنچه تغییر می‌کند، نوع مدل و معیار ارزیابی است.

\section{رگرسیون خطی ساده و چندگانه}

\dmterm{\textbf{رگرسیون خطی ساده}}{Simple Linear Regression} رابطهٔ میان
یک متغیر مستقل $x$ و یک متغیر هدف $y$ را با یک خط راست مدل می‌کند:
\[
\hat{y} = \beta_0 + \beta_1 x
\]
که در آن $\beta_0$ عرض از مبدأ و $\beta_1$ شیب خط است. این ضرایب معمولاً
با روش \dmterm{\textbf{کمترین مربعات}}{Ordinary Least Squares (OLS)}
تخمین زده می‌شوند؛ یعنی مقادیری برای $\beta_0$ و $\beta_1$ انتخاب
می‌شوند که مجموع مربعات خطای پیش‌بینی (فاصلهٔ عمودی هر نقطهٔ داده تا خط
برازش‌شده) را کمینه کنند.

وقتی بیش از یک ویژگی ورودی داشته باشیم، از
\dmterm{\textbf{رگرسیون خطی چندگانه}}{Multiple Linear Regression}
استفاده می‌شود که رابطه را به‌صورت زیر تعمیم می‌دهد:
\[
\hat{y} = \beta_0 + \beta_1 x_1 + \beta_2 x_2 + \cdots + \beta_n x_n
\]

نکتهٔ مهم در رگرسیون چندگانه، توجه به
\dmterm{\textbf{هم‌خطی چندگانه}}{Multicollinearity} است: اگر دو یا چند
ویژگی ورودی به‌شدت با یکدیگر همبسته باشند، تفسیر ضرایب مدل دشوار و
ناپایدار می‌شود، هرچند ممکن است دقت پیش‌بینی کلی مدل تحت تأثیر جدی قرار
نگیرد.

\section{رگرسیون لجستیک}

با وجود نامش، \dmterm{\textbf{رگرسیون لجستیک}}{Logistic Regression}
عملاً یک الگوریتم \textbf{طبقه‌بندی} است، نه رگرسیون به معنای پیش‌بینی
یک مقدار پیوسته؛ این عنوان از تاریخچهٔ توسعهٔ روش باقی مانده است. رگرسیون
لجستیک برای مسائل طبقه‌بندی دوکلاسه به کار می‌رود و به‌جای پیش‌بینی
مستقیم کلاس، \textbf{احتمال} تعلق به کلاس مثبت را تخمین می‌زند.

ایدهٔ اصلی این است که به‌جای برازش یک خط راست (که می‌تواند مقادیری خارج
از بازهٔ $[0, 1]$ تولید کند و برای احتمال بی‌معنی است)، از
\dmterm{\textbf{تابع سیگموید}}{Sigmoid Function} برای نگاشت خروجی مدل
خطی به بازهٔ $[0, 1]$ استفاده می‌شود:
\[
P(y=1 \mid x) = \frac{1}{1 + e^{-(\beta_0 + \beta_1 x_1 + \cdots + \beta_n x_n)}}
\]

در نهایت، با مقایسهٔ این احتمال با یک آستانه (معمولاً ۰.۵)، تصمیم نهایی
کلاس گرفته می‌شود. به همین دلیل، رگرسیون لجستیک را می‌توان پلی میان
مباحث این فصل و مباحث طبقه‌بندی فصل‌های قبل دانست، و از تمام معیارهای
ارزیابی طبقه‌بندی (ماتریس درهم‌ریختگی، دقت مثبت، فراخوانی، منحنی
\LR{ROC}) که در فصل \ref{ch:classification-advanced} دیدیم، برای آن نیز
استفاده می‌شود.

\section{ارزیابی مدل‌های رگرسیون}

از آنجا که خروجی رگرسیون عددی پیوسته است، معیارهای مبتنی بر ماتریس
درهم‌ریختگی (که برای کلاس‌های گسسته تعریف شده‌اند) در اینجا کاربرد
ندارند. به‌جای آن، معیارهای زیر بر پایهٔ فاصلهٔ میان مقدار پیش‌بینی‌شده
$\hat{y}_i$ و مقدار واقعی $y_i$ تعریف می‌شوند.

\paragraph{میانگین مربعات خطا}
\dmterm{\textbf{میانگین مربعات خطا}}{Mean Squared Error (MSE)} میانگین
مجذور تفاضل پیش‌بینی و مقدار واقعی است:
\[
\text{MSE} = \frac{1}{n} \sum_{i=1}^{n} (y_i - \hat{y}_i)^2
\]
به دلیل به توان دو رساندن خطا، این معیار به خطاهای بزرگ (نقاط پرت) حساسیت
بیشتری نسبت به خطاهای کوچک نشان می‌دهد.

\paragraph{جذر میانگین مربعات خطا}
\dmterm{\textbf{جذر میانگین مربعات خطا}}{Root Mean Squared Error (RMSE)}
با گرفتن جذر از \LR{MSE} به دست می‌آید:
\[
\text{RMSE} = \sqrt{\text{MSE}}
\]
مزیت اصلی \LR{RMSE} نسبت به \LR{MSE} این است که واحد آن با واحد متغیر
هدف یکسان است (مثلاً اگر هدف قیمت به تومان باشد، \LR{RMSE} نیز بر حسب
تومان قابل تفسیر است)، در حالی که واحد \LR{MSE} مجذور واحد اصلی است.

\paragraph{ضریب تعیین}
\dmterm{\textbf{ضریب تعیین}}{Coefficient of Determination ($R^2$)} نشان
می‌دهد چه نسبتی از واریانس متغیر هدف توسط مدل توضیح داده می‌شود:
\[
R^2 = 1 - \frac{\sum_{i=1}^{n} (y_i - \hat{y}_i)^2}{\sum_{i=1}^{n} (y_i - \bar{y})^2}
\]
که در آن $\bar{y}$ میانگین مقادیر واقعی است. مقدار $R^2$ نزدیک به ۱
نشان‌دهندهٔ برازش خوب مدل و مقدار نزدیک به صفر (یا حتی منفی) نشان‌دهندهٔ
عملکرد ضعیف آن است.

\section*{پیاده‌سازی در پایتون}

پیاده‌سازی \texttt{LinearRegression} و \texttt{LogisticRegression} از
\texttt{scikit-learn}، به‌همراه محاسبهٔ \texttt{mean\_squared\_error} و
\texttt{r2\_score}، در بخش «فصل ۵» پیوست الف (\ref{app:python}) آمده
است.

\section*{خلاصهٔ فصل}

در این فصل با رگرسیون به‌عنوان دومین وظیفهٔ اصلی یادگیری با نظارت — این
بار برای پیش‌بینی مقادیر پیوسته — آشنا شدیم. رگرسیون خطی ساده و چندگانه
را با روش کمترین مربعات بررسی کردیم، رگرسیون لجستیک را به‌عنوان پلی
میان رگرسیون و طبقه‌بندی معرفی کردیم، و در پایان با معیارهای ارزیابی
مخصوص مسائل رگرسیون ($\text{MSE}$، $\text{RMSE}$، $R^2$) آشنا شدیم. با
این فصل، بخش یادگیری با نظارت این کتاب به پایان می‌رسد. در فصل بعد، به
دنیای یادگیری بدون نظارت و اولین موضوع آن، یعنی قوانین انجمنی، وارد
می‌شویم.

\section*{تمرین‌ها}

\begin{enumerate}
    \item توضیح دهید چرا رگرسیون لجستیک، با وجود نامش، در دستهٔ
          الگوریتم‌های طبقه‌بندی قرار می‌گیرد نه رگرسیون.
    \item فرض کنید یک مدل رگرسیون خطی برای پیش‌بینی قیمت خانه دارید که
          دو ویژگی ورودی «متراژ» و «تعداد اتاق» به‌شدت با هم همبسته‌اند.
          مشکل احتمالی این وضعیت را توضیح دهید.
    \item چرا معیار \LR{RMSE} برای گزارش نتیجه به یک ذی‌نفع غیرفنی
          معمولاً مناسب‌تر از \LR{MSE} است؟
    \item اگر یک مدل رگرسیون مقدار $R^2 = 0.85$ داشته باشد، این عدد را
          به زبان ساده برای یک مدیر کسب‌وکار توضیح دهید.
    \item تابع سیگموید را رسم کنید (روی کاغذ یا با پایتون) و توضیح دهید
          چرا برای تبدیل خروجی یک مدل خطی به احتمال مناسب است.
\end{enumerate}
